\section{Longest Increasing Subsequence}

\subsection*{Framework of Dynamic Programming Algorithm}

A typical dynamic programming algorithm consists of following steps:
\vspace*{-\topsep}
\begin{enumerate}
\item Define subproblems, %~(or equivalently, partition the search space into sub-spaces),
$P_1, P_2, \cdots, P_n$. Solving all these subproblems will solve the original problem.
sometimes the original problem is $P_n$, but sometimes it is not.
\item Develop a recurrence for these subproblems. Formally, let $opt(P_k)$ be the optimal
solution of $P_k$. A recurrence is a function that calculates
$opt(P_{k+1})$ using $opt(P_1), opt(P_2), \cdots, opt(P_k)$.
\item Fill a \emph{dynamic programming table},
in which each entry stores the optimal solution of a subproblem, 
using above recurrence.
\end{enumerate}

\subsection*{Optimal Substrcture Property}

Let $P_1, P_2, P_3, \cdots, P_n$ be a series of subproblems defined for original problem $P$.
If  $opt(P_k)$ can be calculated solely from the \emph{optimal solution} of previous subproblems,
i.e., $opt(P_k) = f(opt(P_1), opt(P_2), \cdots, opt(P_{k-1}))$,
then we say $P$ and its subproblems $\{P_1, P_2, \cdots, P_n\}$ satisfies the optimal substructure property.
Intuitively, the optimal substructure property states that
the optimal solution of a larger problem can be \emph{assembled}
using the optimal solutions of the smaller subproblems.

How to determine if a definition of subproblems satisfies the optimal substructure property?
Usually we can verify if its opporsite side is true, i.e., if the optimal solution
of a bigger subproblem \emph{implies} the optimal solution of
the smaller subproblems.
%Take the shortest path problem as an example we showed before:
%Let $p = s \to w \to \cdots \to u \to v$
%be the shortest path from $s$ to $v$.
%Then clearly its sub-path $s\to w \to \cdots \to u$ is the shortest path from
%$s$ to $u$~(because, if there exists a shorter path from $s$ to $u$
%then we can find a shorter one from $u$ to $v$ than $p$).
%This suggests that the shortest path problem
%satisfies the optimal substructure property.
We will use examples to demonstrate it.
%and how to determine a proper definition accordingly.


\subsection*{Longest Increasing Subsequences}

Let $A[1\cdots n]$ be an array with $n$ distinct integers.
A \emph{subsequence} of array $A$ is a subset of $A$ but follows
the same order. For example, $(A[2], A[4], A[5])$ is a subsequence of $A$
while $(A[3], A[1], A[4])$ is not.
An \emph{increasing subsequence}~(IS) of $A$ means the values of the
subsequence are increasing.
Formally, an IS of length $k$ can be written
as $(A[i_1], A[i_2], \cdots, A[i_k])$, where $1\le i_1 < i_2 < \cdots < i_k \le n$,
and $A[i_1] < A[i_2] < \cdots < A[i_k]$.
The \emph{longest increasing subsequence problem}~(LIS problem) is, given array $A[1\cdots n]$,
to seek an IS with maximum length.
(We assume that $A$ contains $n$ \emph{distinct} numbers for sake
of simplicity; all algorithms introduced later can be modified to handle duplicates.)

For example, let $A = [1, 8, 3, 9, 4, 7, 5]$. Then one LIS is $[1, 3, 4, 7]$; %The optimal solution may not be unique though. 
another LIS is $[1, 3, 4, 5]$.

We want to design a dynamic programming algorithm to solve this problem.
The first step is to define subproblems. Naturally, we define $P_k$
as to seek the LIS for $A[1\cdots k]$. Solving $P_n$ gives the LIS
of $A[1\cdots n]$. Is this a good definition? Let's verify if substructures
of $opt(P_k)$ implies the optimal solutions of smaller subproblems.
Consider this example:

\begin{displaymath}
\begin{array}{rllllllll}
A & = & [3, &8, &9, &4, &6, &7, &5] \\
opt(P_1) &=& [3] \\
opt(P_2) &=& [3, &8] \\
opt(P_3) &=& [3, &8, &9] \\
opt(P_4) &=& [3, &8, &9] \\
opt(P_5) &=& [3, &&&4, &6] &&\quad (\textrm{or } [3, 8, 9]) \\
opt(P_6) &=& [3, &&&4, &6, &7] \\
opt(P_7) &=& [3, &&&4, &6, &7]
\end{array}
\end{displaymath}

Consider $opt(P_6) = [3, 4, 6, 7]$ and its 3 substructures $[3]$, $[3, 4]$, and $[3, 4, 6]$.
The first and the last one are indeed the optimal solutions of $P_1$ and $P_5$,
but $[3, 4]$ is not the optimal solution of $P_4$.
This suggests that these subproblems do not satisfy the optimal substructure property and
not suitable for designing a DP algorithm.

Let's have another try. Define $P_k$ as to find the LIS of $A[1\cdots k]$ that ends at $A[k]$.
This is a common technique that introduce a constraint to the definition of the subproblems.
Here we require that $opt(P_k)$ must use, and end at, $A[k]$. Let's manually work out
the optimal solutions, following this new defintion, on above example:
\begin{displaymath}
\begin{array}{rllllllll}
A & = & [3, &8, &9, &4, &6, &7, &5] \\
opt(P_1) &=& [3] \\
opt(P_2) &=& [3, &8] \\
opt(P_3) &=& [3, &8, &9] \\
opt(P_4) &=& [3, &&&4] \\
opt(P_5) &=& [3, &&&4, &6] \\
opt(P_6) &=& [3, &&&4, &6, &7] \\
opt(P_7) &=& [3, &&&4, &&&5] \\
\end{array}
\end{displaymath}

You may verify that, any substructure of any optimal solution implies the optimal
solution of the smaller subproblem. In fact, we can formally prove it.

%Recall that, in order to give a proper definition of subproblems,
%we can try ``reverse engineering'', i.e., starting from the optimal
%solution, trying to determine if it \emph{implies} the optimal solution
%of certain subproblem. If it does, then that subproblem is likely the one
%we could use in designing the DP algorithm.
%
%Let's try this approach with on above example. We know that $[1, 3, 4, 7]$
%is one optimal solution. Does its substructure, $[1, 3, 4]$, give the optimal
%solution of certain subproblem? A natural guess would be that, $[1, 3, 4]$,
%is the LIS of $A[1\cdots 6] = [1, 5, 3, 8, 9, 4]$.
%But, a further look denies this conjecture, as $[1, 3, 8, 9]$ is a longer
%IS of $A[1\cdots 6]$.
%Let's have another try. We still consider $A[1\cdots 6]$, but require that
%the last number of the IS must be $A[6] = 4$.
%(Note: this is again a commonly-used technique, i.e., introducing a variable or a constraint
%into the definition of the subproblems.)
%Among all such ISs of $A[1\cdots 6]$, is $[1, 3, 4]$ the optimal one?
%Yes! Why? We can prove by contradiction: suppose that there is a longer IS of $A[1\cdots 6]$
%that ends at $A[6] = 4$, we can concatenate it with $A[7] = 7$ to obtain 
%a longer, than $[1, 3, 4, 7]$, IS of $A$, contradicting
%to the optimality of $[1, 3, 4, 7]$.

%We summarize above example with the following optimal substructure statement.
\begin{fact}\label{lis-fact1}
Let $(A[i_1], A[i_2], \cdots, A[i_k])$ be one LIS of $A[1\cdots i_k]$.
Then  $(A[i_1], A[i_2], \cdots, A[i_j])$ must be one LIS of $A[1\cdots i_j]$
for any $j$, $j=1, 2, \cdots, k-1$.
\end{fact}
\emph{Proof.} We prove it by contradiction. Assume that, for some $j$, 
$(A[i_1], A[i_2], \cdots, A[i_j])$ is not optimal, i.e.,
there exists an IS ending at $A[i_j]$ whose length is greater than $j$.
Extending this IS by appending $(A[i_{j+1}], A[i_{j+2}], \cdots, A[i_k])$
gives an IS, because $A[i_j] < A[i_{j+1}]$, whose size is larger than $k$.
This contradicts to the optimality of $(A[i_1], A[i_2], \cdots, A[i_k])$. \qed

We proved that above definition of subproblems $\{P_k\}$ satisfy the optimal substructur property.
We proceed with developing the recurrence.
%Above property directly suggests the definition of subproblems.
%We define $opt(k)$ as the LIS of $A[1\cdots k]$ with $A[k]$ being the last number,
%and define $L(k)$ as the length of $opt(k)$.
%In other words, among all ISs of $A[1\cdots k]$ with last number being $A[k]$,
%$opt(k)$ is the optimal solution and $L(k)$ is its length.
%
%We defined $n$ subproblems. Suppose we can solve all of them.
%What's the LIS of $A$?
%Clearly, that's $\max_{1\le k \le n} L(k)$.
%
%We now develop a recurrence. 
How to construct $opt(P_k)$ given $opt(P_j)$, $j = 1, 2, \cdots, k-1$?
A useful technique is to enumerate the last step.
%we enumerate the last step. By definition of the subproblem, $A[k]$ is the last number of $opt(k)$.
For the LIS problem, it is to enumerate the second last number~(the last number of $opt(P_k)$ is $A[k]$ by its definition).
The second last number must be one of $A[j]$ satisfying $1\le j < k$ and $A[j] < A[k]$.
%We don't know which one is correct so we enumerate all possibilities.  (And once we enumerate, we can assume we know it.)
Suppose that $A[j]$ is the second last number of $opt(P_k)$, then according to Fact~\ref{lis-fact1},
$opt(P_k)$ can be obtained by concatenating $opt(P_j)$ and $A[k]$.
Hence, seeking $opt(P_k)$ is equivalent to finding $j$ such that $opt(P_j)$ is the longest. 
Let $L(k)$ to denote the length of $opt(P_k)$. We can write the recurrence:
$$\textstyle L(k) = 1 + \max_{1 \le j < k, A[j] < A[k]} L[j].$$

